\let\negmedspace\undefined
\let\negthickspace\undefined
\documentclass[journal]{IEEEtran}
\usepackage[a5paper, margin=10mm, onecolumn]{geometry}
%\usepackage{lmodern} % Ensure lmodern is loaded for pdflatex
\usepackage{tfrupee} % Include tfrupee package

\setlength{\headheight}{1cm} % Set the height of the header box
\setlength{\headsep}{0mm}     % Set the distance between the header box and the top of the text

\usepackage{gvv-book}
\usepackage{gvv}
\usepackage{cite}
\usepackage{amsmath,amssymb,amsfonts,amsthm}
\usepackage{algorithmic}
\usepackage{graphicx}
\usepackage{textcomp}
\usepackage{xcolor}
%\usepackage{txfonts}
\usepackage{listings}
\usepackage{enumitem}
\usepackage{mathtools}
\usepackage{gensymb}
\usepackage{comment}
\usepackage[breaklinks=true]{hyperref}
\usepackage{tkz-euclide} 
\usepackage{listings}
% \usepackage{gvv}                                        
\def\inputGnumericTable{}                                 
\usepackage[latin1]{inputenc}                                
\usepackage{color}                                            
\usepackage{array}                                            
\usepackage{longtable}                                       
\usepackage{calc}                                             
\usepackage{multirow}                                         
\usepackage{hhline}                                           
\usepackage{ifthen}                                           
\usepackage{lscape}
\usepackage{circuitikz}
\tikzstyle{block} = [rectangle, draw, fill=blue!20, 
    text width=4em, text centered, rounded corners, minimum height=3em]
\tikzstyle{sum} = [draw, fill=blue!10, circle, minimum size=1cm, node distance=1.5cm]
\tikzstyle{input} = [coordinate]
\tikzstyle{output} = [coordinate]


\begin{document}

\bibliographystyle{IEEEtran}
\vspace{3cm}

\title{8.3.10}
\author{AI25BTECH11030 -Sarvesh Tamgade}
{\let\newpage\relax\maketitle}

\renewcommand{\thefigure}{\theenumi}
\renewcommand{\thetable}{\theenumi}
\setlength{\intextsep}{10pt} 


\numberwithin{equation}{enumi}
\numberwithin{figure}{enumi}
\renewcommand{\thetable}{\theenumi}


\textbf{Question}:
If the eccentricity of an ellipse is $\frac{5}{8}$ and the distance between its foci is $10$, then find the latus rectum of the ellipse.


\textbf{Solution}:  
Given:
\begin{itemize}
\item Eccentricity: $e = \frac{5}{8}$
\item Distance between foci: $D = 10$
\end{itemize}

\textbf{Step 1: Express the foci for standard ellipse}

From B.2.7.3, the foci for standard ellipse are at:

\begin{align}
\vec{F} = \pm\sqrt{\frac{f_0}{\lambda_2(1-e^2)}} \vec{e_1}
\end{align}

where $f_0$ is the scalar term in the conic equation and $\lambda_2$ is an eigenvalue of matrix $\vec{V}$.

\textbf{Step 2: Distance between foci}

The distance between foci is:

\begin{align}
D = 2\sqrt{\frac{f_0}{\lambda_2(1-e^2)}}
\end{align}

\textbf{Step 3: Solve for $f_0$}

From equation (2):

\begin{align}
\frac{D}{2} &= \sqrt{\frac{f_0}{\lambda_2(1-e^2)}}
\end{align}

Squaring both sides:

\begin{align}
\frac{D^2}{4} &= \frac{f_0}{\lambda_2(1-e^2)}
\end{align}

Therefore:

\begin{align}
f_0 &= \frac{D^2 \lambda_2(1-e^2)}{4}
\end{align}

\textbf{Step 4: Eigenvalues for standard ellipse}

For the standard ellipse, from the conic section theory:

\begin{align}
\lambda_1 &= \frac{||\vec{n}||^2}{1-e^2}
\end{align}

\begin{align}
\lambda_2 &= ||\vec{n}||^2
\end{align}

For the standard form with $||\vec{n}||^2 = 1$:

\begin{align}
\lambda_1 &= \frac{1}{1-e^2}
\end{align}

\begin{align}
\lambda_2 &= 1
\end{align}

\textbf{Step 5: Calculate $f_0$ with $\lambda_2 = 1$}

Substituting $\lambda_2 = 1$ into equation (5):

\begin{align}
f_0 &= \frac{D^2(1-e^2)}{4}
\end{align}

Therefore:

\begin{align}
\sqrt{f_0} &= \frac{D\sqrt{1-e^2}}{2}
\end{align}

\textbf{Step 6: Matrix formulation for latus rectum}
The length of latus rectum is given by :

\begin{align}
\ell = \frac{2\sqrt{f_0 \lambda_1}}{\lambda_2}
\end{align}

Substituting the values of $f_0$, $\lambda_1$ and $\lambda_2$ 
\begin{align}
\ell &= \frac{D(1-e^2)}{e}
\end{align}

\textbf{Step 7: Substitute the given values}

Given $e = \frac{5}{8}$ and $D = 10$:

\begin{align}
1 - e^2 &= 1 - \left(\frac{5}{8}\right)^2 = 1 - \frac{25}{64} = \frac{39}{64}
\end{align}

Substituting into the latus rectum formula:
\begin{align}
\ell &= \frac{390 \times 8}{64 \times 5} = \frac{3120}{320} = \frac{39}{4}
\end{align}

\textbf{Final Answer:}
The length of latus rectum is :

\begin{align}
\boxed{\ell = \frac{39}{4}}
\end{align}
\newpage
\begin{figure}[htbp]
    \centering
    \includegraphics[width=0.65\linewidth]{FIG/fig1.png}
    \caption{Vector Representation}
    \label{fig:FIG/fig1.png}
    \end{figure}

\end{document}  